\section{Companion lab: a shortest route through a maze}
\label{revision-lab:maze}

Sorting required a useful order for array entries. A search also needs an order
for unfinished work. In this lab, an array of cell IDs records places waiting to
be explored. Taking the oldest waiting cell first lets the search expand outward
in layers of increasing distance. The result is a shortest route through an
unweighted maze, together with enough intermediate state to explain that route.

This is original project \texttt{010-maze-explorer}. Its contract is
\texttt{assignment.md} in that package. The revised work area is
\texttt{maze-search.asm}, under \texttt{revision/labs/files}. Both public checks
remain in the original package's \texttt{checks} directory. Implement the search and its
\texttt{render} function. The starter's \texttt{break} instructions deliberately
stop execution until those missing bodies are replaced. Organize private functions
and registers as you choose; the public allocations, observations and renderer
interface establish what another program can inspect.

You will use byte arrays, row-major addresses, loop invariants and calls that
preserve live state. The new ideas are graph adjacency, discovery-time marking,
a first-in, first-out frontier, predecessor links and a shortest-distance argument.
Breadth-first search, or BFS, is the conventional name for this layer-by-layer
search.\footnote{For the layer argument and predecessor reconstruction, see MIT
OpenCourseWare, \emph{6.006 Introduction to Algorithms}, Spring 2020, Recitation 9,
pp.~3--4: \url{https://ocw.mit.edu/courses/6-006-introduction-to-algorithms-spring-2020/resources/mit6_006s20_r09/}.}

\subsection*{Represent cells and legal moves separately}

The input is an eight-by-eight array of bytes. A zero \texttt{walls} byte denotes
an open cell and one denotes a wall. The word inputs \texttt{start} and
\texttt{goal} name open cells by IDs from zero through 63. A cell at column $x$
and row $y$ has ID $8y+x$. Its byte entry is at $walls+ID$; its pixel word is at
$frame+4ID$. An ID is neither of those addresses.

Treat each open cell as a vertex of a graph, and each legal one-cell move as an
edge. Moves are north, east, south or west, with cost one each. Diagonals are
excluded. The graph is implicit: the program computes possible neighbors from
coordinates rather than storing a linked collection for every cell.

Derive $x=ID\bmod8$ and $y=\lfloor ID/8\rfloor$, apply a direction displacement,
then reject a coordinate outside zero through seven. Only after these checks
should the candidate be flattened back to an ID. Cell seven's eastern candidate
has column eight and must be rejected. Adding one to its ID gives eight, which
is a valid array index but belongs to the next row. An array bound alone cannot
establish adjacency.

The complete \texttt{maze-neighbor.asm} experiment isolates this geometry. Its
\texttt{neighbor\_id} function takes a valid cell ID in \texttt{\$a0} and a
direction index in \texttt{\$a1}: zero north, one east, two south, three west.
It returns the neighbor ID in \texttt{\$v0}, or $-1$ when the move leaves the grid.
It does not inspect walls; an in-bounds cell can still be blocked. The direction
tables contain signed bytes, so \texttt{lb} must preserve the $-1$ displacements.

\noindent\begin{minipage}{\linewidth}
\MARSCodeLabel{Check both coordinates before producing an ID}
\begin{minted}{mars-idea}
neighbor_id:
    andi $t0, $a0, 7
    srl $t1, $a0, 3
    la $t2, dx
    daddu $t2, $t2, $a1
    lb $t3, 0($t2)
    addu $t0, $t0, $t3
    bltzc $t0, neighbor_absent
    nop
    li $t4, 8
    bgec $t0, $t4, neighbor_absent
    nop
    la $t2, dy
    daddu $t2, $t2, $a1
    lb $t3, 0($t2)
    addu $t1, $t1, $t3
    bltzc $t1, neighbor_absent
    nop
    bgec $t1, $t4, neighbor_absent
    nop
    sll $v0, $t1, 3
    addu $v0, $v0, $t0
    jrc $ra
neighbor_absent:
    li $v0, -1
    jrc $ra
\end{minted}
\end{minipage}

The bit mask extracts the low three column bits, and the shift extracts the row.
After applying a displacement, the two signed checks on each coordinate reject
negative and too-large results. Pointer additions use full-width instructions;
the coordinates themselves fit comfortably in signed words. Saved registers and
the stack pointer are untouched. This helper is an experiment you may call or
adapt, not a required private decomposition of the full search.

Run \texttt{revision/labs/checks/maze-neighbor.kts} against that complete source
with the CLI. It checks all 256 valid cell/direction combinations and the supplied
caller. Before running it, predict east from seven, west from eight, north from
zero, and south from 55. Explain each answer with coordinates. In a copy, replace
\texttt{lb} with \texttt{lbu} and find which directions now fail. Restore the
signed loads before using the experiment as a search component.

\subsection*{Make discovery a single event}

Five 64-byte arrays hold different facts. \texttt{states} contains zero for unseen,
one for discovered and waiting, and two for processed. \texttt{parents} stores the
ID of the cell that first discovered an entry. \texttt{distance} stores its number
of edges from the start. \texttt{queue} retains discovery IDs in FIFO order, and
\texttt{path} eventually marks the chosen route. Unknown parents and distances
are 255; queue and path bytes begin at zero. Initialize every entry on every launch,
including those that the eventual search will never reach.

These arrays are byte-indexed, so load IDs and distances with \texttt{lbu}.
The sentinel 255 is not a valid ID and is not a distance to add one to. A simple
route in a graph of at most 64 cells has at most 63 edges, leaving the sentinel
outside the valid distance range. Word metadata uses \texttt{lw}/\texttt{sw};
all public byte-array bases and metadata words are four-byte aligned. Preserve
the declared guard words and execute no stores to \texttt{walls}, \texttt{start}
or \texttt{goal}, even if later writes could restore them.

Initialize the start as waiting with distance zero and no parent, then append
its ID. For each newly discovered neighbor, set its waiting state, parent and
distance before appending it. The parent is the current cell; the distance is
that current cell's known distance plus one. Skip walls and every already
discovered cell, including waiting cells that have not yet been removed.

The four-cell room shows why the mark belongs at discovery. Zero discovers both
one and eight.
Processing one can discover nine. When eight later examines nine, it must see the
existing mark and leave nine's parent and queue entry alone. If marking waited
until dequeue, both predecessors could append nine before either copy was
processed. Cycles would repeat work, invalidate unique history and undermine the
queue capacity argument.

\subsection*{Use the queue to explain the shortest distance}

The monotone queue uses two word indices. \texttt{head} names its next unread
byte and \texttt{tail} its first unused byte. Pending work occupies the half-open
interval $[head,tail)$; retained history occupies $[0,tail)$. The invariant is
$0\leq head\leq tail\leq64$. An empty pending queue has equal indices, even
though earlier queue entries still contain IDs. Do not clear a removed entry.
The unused suffix remains zero.

This queue does not wrap like the preceding ring assignment. Each cell is
appended at most once, so the entire search needs at most 64 slots, regardless
of how many slots remain pending at one moment. Appending ID zero is valid:
use head and tail to identify occupied entries instead of treating a zero byte
as an end marker.

The FIFO order is what gives the distance its meaning. Initially only distance
zero is waiting. When a cell at distance $d$ is removed, its unseen neighbors
receive $d+1$ and go behind work already waiting. All cells at smaller distances
are therefore processed before any newly appended farther layer. If a neighbor
had a shorter route, the preceding cell on that route would have been processed
in an earlier layer and discovered it already. Its first assigned distance is
thus shortest. This argument depends on equal edge costs and FIFO removal;
a stack or weighted terrain would require a different argument.

\begin{figure}[htbp]
  \centering
  \begin{tikzpicture}[x=17mm,y=-17mm,alt={The upper-left two by two room contains cell zero at top left, one at top right, eight at bottom left and nine at bottom right. All other maze cells are walls. Orthogonal links form a cycle. Two routes from zero to nine have two edges: through one or through eight. Arrows highlight the first route; a parent walk traverses it backward from nine to one to zero.}]
    \draw[thick] (0,0) rectangle (2,2);
    \draw (1,0)--(1,2) (0,1)--(2,1);
    \node at (.5,.5) {0}; \node at (1.5,.5) {1};
    \node at (.5,1.5) {8}; \node at (1.5,1.5) {9};
    \draw[->,very thick] (.72,.5)--(1.28,.5);
    \draw[->,very thick] (1.5,.72)--(1.5,1.28);
    \node[anchor=west,align=left,text width=62mm] at (2.4,.55) {Two equal-length routes:\\$0\to1\to9$\\$0\to8\to9$};
    \node[anchor=west,align=left,text width=62mm] at (2.4,1.5) {For the highlighted route,\\parents lead backward:\\$9\to1\to0$.};
  \end{tikzpicture}
  \caption{Cell IDs use the full maze's eight-cell row stride, even in this small open room. A route chooses edges; a parent entry records the preceding cell's ID.}
  \label{fig:revision-maze-room}
\end{figure}

\begin{table}[htbp]
  \centering
  \small
  \begin{tabular}{lrrlll}
    \toprule
    Boundary & Head & Tail & Queue history    & Waiting & New parent           \\
    \midrule
    Initial  & 0    & 1    & \texttt{0}       & 0       & none                 \\
    After 0  & 1    & 3    & \texttt{0,1,8}   & 1,8     & $1\gets0$, $8\gets0$ \\
    After 1  & 2    & 4    & \texttt{0,1,8,9} & 8,9     & $9\gets1$            \\
    After 8  & 3    & 4    & \texttt{0,1,8,9} & 9       & none                 \\
    After 9  & 4    & 4    & \texttt{0,1,8,9} & none    & none                 \\
    \bottomrule
  \end{tabular}
  \caption{A FIFO trace in the four-cell room, using north, east, south, west neighbor order. Head names the next unread queue position; tail is the first unused position. History is retained after removal. The final path is $0,1,9$, with two edges.}
  \label{tab:revision-maze-frontier}
\end{table}


Follow Table~\ref{tab:revision-maze-frontier} by hand before implementing the
loop. With north, east, south, west order, nine first receives parent one.
With west, south, east, north order, the waiting order after zero is eight, one,
and nine first receives parent eight. Both routes have two edges. The public
checks accept either order and follow your actual queue history; they require a
shortest valid chain, not the reference program's tie choice.

Remove one pending ID, advance head, and mark that current cell processed.
If it is the goal, skip its neighbor expansion. Otherwise inspect all four
candidates, appending only unseen open neighbors. Publish after the complete
expansion, then choose the next item. A reachable goal is finalized when dequeued
under this contract, even if its shortest distance became known earlier.

\subsection*{Publish a state that another observer can interpret}

Set phase zero after initialization, phase one after each complete dequeue and
expansion, and phase two after success or exhaustion. At every publication,
\texttt{head}, \texttt{tail} and \texttt{current} describe that boundary. Initially
current is $-1$. Result and length remain $-1$ throughout phases zero and one,
including the phase-one publication of a dequeued goal. The path mask remains
zero until reconstruction. A cell in the middle of expansion can have processed
state internally, but phase-one observers must receive all its discoveries together.

The leaf-style \texttt{render} interface reads \texttt{walls}, \texttt{states}
and \texttt{path}, writes all 64 pixel words at \texttt{frame}, and changes no
input array or metadata. It preserves full saved-register values and the stack
pointer. It reaches \texttt{maze\_ready} exactly once with the frame complete
and preserved state restored, then returns through the supplied return address.
The \texttt{done} label supports an independent call. If you introduce nested
calls or temporary stack storage, apply the earlier call-frame rules.

Color priority is wall, path, then state. In hexadecimal RGB, walls are
\texttt{202020}, path cells \texttt{0072B2}, unseen cells \texttt{FFFFFF}, waiting
cells \texttt{E69F00}, and processed cells \texttt{009E73}. Repaint every pixel,
including pixels whose new color is white or dark. The colors help locate state;
array values and IDs provide the explanation without depending on color perception.
The framebuffer is the first static allocation at hexadecimal \texttt{10010000};
its immediate guard is at \texttt{10010100}. Other named arrays may move.

Set a breakpoint at \texttt{maze\_ready} in Debug. Inspect phase zero and the
single waiting start, then continue through the four-cell trace. At each phase-one
pause, derive the pending range from head and tail and inspect the newly appended
IDs and their parents. Step beyond the current occurrence before waiting for the
next one. A fast Run may show only the final frame; screen refreshes are not a
count of completed expansions.

\subsection*{Follow parents and distinguish failure from a zero-length route}

On success, set result to one and length to the goal's distance. Starting at the
goal, mark that cell in \texttt{path}, then follow its parent until reaching the
start. Mark the start too, and stop before reading its sentinel parent. Each
parent must be an adjacent open cell whose distance is one smaller. This decreasing
nonnegative distance explains why the chain terminates without a cycle.

For route $0,1,9$, reconstruct backward as $9,1,0$. The mask has three marked
cells but length is two: the contract counts edges. If start equals goal, the
search still initializes and dequeues that cell, then succeeds with length zero
and one marked cell. It does not expand the goal's neighbors.

If pending work becomes empty without removing the goal, set result to zero,
retain length $-1$ and leave the path mask all zero. Reach phase two, redraw and
exit normally in either case. An unreachable goal and a zero-edge success have
different result, length and path states; do not infer either from an empty queue
alone.

Figure~\ref{fig:revision-maze-routes} compares the two valid reconstructed chains.
Use the IDs and publication outlines as well as the colors when comparing your
result. A different parent for nine is acceptable when its chain remains valid
and shortest.

\begin{figure}[htbp]
\centering
\begin{tikzpicture}[alt={Two actual native Bitmap Display observations cropped to the upper-left two by two open room. The left route marks cells zero, one and nine; cell eight is processed but off the route. The right route marks zero, eight and nine; cell one is processed but off the route. Native publication numbers identify all four cells and double outlines identify the three route cells in each panel, independently of color. Both routes have two edges.}]
\node[anchor=north west,inner sep=0,font=\small\bfseries] (lefttitle) at (0,0) {1. East before south};
\node[anchor=north west,inner sep=0] (left) at ([yshift=-2mm]lefttitle.south west) {\MARSRevisionCrop{revision/captures/maze/maze-east-first.png}{0.28\linewidth}{0.203906250}{0.698333333}{0.648046875}{0.091111111}};
\coordinate (rightorigin) at ([xshift=12mm]left.north east);
\node[anchor=north west,inner sep=0,font=\small\bfseries] (righttitle) at (rightorigin |- lefttitle.north west) {2. South before east};
\node[anchor=north west,inner sep=0] (right) at (righttitle.south west |- left.north west) {\MARSRevisionCrop{revision/captures/maze/maze-south-first.png}{0.28\linewidth}{0.203906250}{0.698333333}{0.648046875}{0.091111111}};
\foreach \panel in {left,right} {
 \begin{scope}[shift={(\panel.south west)},x={($(\panel.south east)-(\panel.south west)$)},y={($(\panel.north west)-(\panel.south west)$)}]
 \draw[black,line width=.4pt] (0,0) rectangle (1,1) (.5,0)--(.5,1) (0,.5)--(1,.5);
 \foreach \x/\y/\id in {.25/.75/0,.75/.75/1,.25/.25/8,.75/.25/9} {\node[font=\normalsize,fill=white,inner sep=1.5pt] at (\x,\y) {\id};}
 \end{scope}
}
\foreach \panel/\x/\y in {left/.02/.52,left/.52/.52,left/.52/.02,right/.02/.52,right/.02/.02,right/.52/.02} {
 \begin{scope}[shift={(\panel.south west)},x={($(\panel.south east)-(\panel.south west)$)},y={($(\panel.north west)-(\panel.south west)$)}]
 \draw[black,line width=1.8pt] (\x,\y) rectangle ++(.46,.46);
 \draw[white,line width=.7pt] (\x,\y) rectangle ++(.46,.46);
 \end{scope}
}
\end{tikzpicture}
\caption{Actual light-theme Debug observations of the open room; other cells are walls. Native LaTeX numbers identify cells, and double outlines mark each route. Neighbor order changes the chosen shortest chain; both edge counts are two.}
\label{fig:revision-maze-routes}
\end{figure}


In a copy, make only cells seven and eight open, with seven the start and eight
the goal. A result of one reveals row wrapping. Use cells zero and nine for a
diagonal-only case; they are disconnected too. For the original four-cell room,
delay the discovery mark in a disposable implementation and inspect the first
duplicate append. Repair that transition before trying a larger open maze.

\subsection*{Build and check the search independently}

Begin with full initialization and the phase-zero frame. Then implement one
dequeue and its complete expansion; use the table to inspect the first two such
boundaries. Generalize that loop, add exhaustion and success, and finally reconstruct
the predecessor chain. These stages separate faulty geometry, duplicate discovery
and incorrect reconstruction instead of hiding all three behind one final picture.
Keep the renderer independently callable so its faults can be investigated without
running a complete search.

Run both package checks against your revised source:
\begin{verbatim}
java -jar /path/to/cli.jar grade \
  packages/010-maze-explorer/checks/public.kts \
  /path/to/your/maze-search.asm
java -jar /path/to/cli.jar grade \
  packages/010-maze-explorer/checks/render.kts \
  /path/to/your/maze-search.asm
\end{verbatim}
The checks inspect every published transition and accept your scratch-register
choices. Renderer checks supply independent arrays and saved registers. Keep a
failing input and locate its first disagreement: a missing neighbor, duplicate
append or changed parent can precede the wrong final mask. Repair that transition
and rerun both checks, rather than patching only the picture.

After the basic search passes, change the open room to force a detour and justify
its new edge count. Predict which parent changes if you reverse neighbor order.
Then measure the largest pending queue size at completed expansion boundaries;
state that sampling boundary, since it excludes the internal moment after a
removal but before its new appends. For a larger extension, design a rectangular
interface with explicit row stride and capacity. Do not carry the eight-column
mask or 64-slot proof into that interface without deriving their replacements.
